\begin{example}
Soient \todo{Lecture 3} deux polynomes $f(x)=x^{2}+x \in \mathbf{Z}_{2}[x]$ et $g(x)=x^{3}+x \in \mathbf{Z}_{2}[x]$ on a que $f_{g}=f_{f}=0$ mais $f(x)\neq g(x)$.
\end{example}
\section{L'interpolation}

\begin{theorem}
Soit $K$ un corps et $r_{0},\dots,r_{n}\in K$ distincts (c-a-d $r_{i}\neq r_{j}$ pour $i\neq j$) et $y_{0},\dots,y_{n}\in K$. Il existe un unique polynome $f(x)=a_{0}+a_{1}x+\dots+a_{n}x^{n}$ de degre au plus $n$ tel que :
$$
\forall i=0,\dots,n:f(r_{i})=y_{i}
$$

\end{theorem}
\begin{proof}
On commence par observer que pour tout $i$ on a que $f(r_{i})=y_{i}$ si et seulement si 
$$
\begin{pmatrix}
1 & r_{0} & r_{0}^{2} & \dots & r_{0}^{n} \\
1 & r_{1} & r_{1}^{2} & \dots & r_{1}^{n} \\
 &  & \vdots &  &  \\
1 & r_{n} & r_{n}^{2} & \dots & r_{n}^{n}
\end{pmatrix}
\begin{bmatrix}
a_{0} \\
\vdots \\
a_{n}
\end{bmatrix}=\begin{pmatrix}
y_{0} \\
\vdots \\
y_{n}
\end{pmatrix}
$$
on appelle la matrice de gauche une matrice de Van der Monde $V(r_{0},\dots,r_{n})\in K^{(n+1)\times(n+1)}$. Il nous reste a montrer que $V(r_{0},\dots,r_{n})\in K^{(n+1)\times (n+1)}$ est inversible (c-a-d $\det(V(r_{0},\dots,r_{n}))\neq 0$). 

On le montre par recurrence sur $n$. On commence par etablir la validité de l'hypothese de recurrence pour $n=0$ qui donne la matrice $V(r_{0})=\begin{pmatrix}1\end{pmatrix}$ alors $\det(V(r_{0}))=1\neq 0$. On admet l'hypothese valide pour $n-1$, on veut calculer
$$
\det \begin{pmatrix}
1 & r_{0} & r_{0}^{2} & \dots & r_{0}^{n-1} & r_{0}^{n} \\
\vdots & \vdots & \vdots & \vdots & \vdots & \vdots  \\
1 & r_{n} & r^{2}_{n} & \dots & r_{n}^{n-1} & r_{n}^{n}
\end{pmatrix}
$$
ou on peut appliquer des operations elementaires sur les colonnes pour obtenir
$$
\det \begin{pmatrix}
1 & 0 & \dots & 0 & 0 \\
1 & (r_{1}-r_{0}) & \dots & (r_{2}-r_{0})r_{1}^{n-2} & (r_{1}-r_{0})r_{1}^{n-1} \\
\vdots &\vdots  &  \vdots &\vdots  &\vdots  \\
1 & (r_{n}-r_{0}) & \dots & (r_{n}-r_{0})r_{n}^{n-2} & (r_{n}-r_{0})r^{n-1}_{n}
\end{pmatrix}
$$
qui est egal a
$$
\det \left[\begin{pmatrix}
(r_{1}-r_{0}) &   &   &  \\
 & (r_2-r_{0})    & &  \\
 & &  \ddots \\

 &  &   &  (r_{n}-r_{0})
\end{pmatrix}\cdot V(r_{1},\dots,r_{n})\right]\neq 0
$$
\end{proof}
\begin{example}
On cherche $f(x)\in \mathbf{Z}_{5}[x]$ de $\deg f\leq3$ tel que 
\begin{align*}
f(0) & =2 \\
f(1) & =2 \\
f(2) & =2 \\
f(3) & =3
\end{align*}
on a 
$$
\begin{pmatrix}
1 & 0 & 0 & 0 \\
1 & 1 & 1 & 1 \\
1 & 2 & 4 & 3 \\
1 & 3 & 4 & 2
\end{pmatrix}\begin{pmatrix}
a_{0} \\
a_{1} \\
a_{2} \\
a_{3}
\end{pmatrix}=\begin{pmatrix}
2 \\
2 \\
2 \\
3
\end{pmatrix}
$$
qui donne une solution unique
$$
f(x)=2+2x+2x^{2}+x^{3}
$$
\end{example}
\begin{corollary}
Soit $K$ un corps infini. Deux polynomes $h(x)$ et $g(x)$ dans $K[x]$ sont egaux si et seulement si les applications $f_{g}$ et $f_{h}$ sont les memes.
\end{corollary}

\section{Division avec reste}

\begin{example}
Dans $\mathbf{Z}_{3}$ on divise $p(x)=x^{5}+2x^{2}+1$ par $q(x)=2x^{3}+x+1$ qui nous donne $s(x)=2x^{2}+2$ avec une reste $r(x)=x+2$ ainsi on a que $p(x)=q(x)s(x)+r(x)$ et $\deg r<\deg g$.
\end{example}
\begin{theorem}
Soit $R$ un anneau et $f,g\in \mathbf{R}[x]$ tel que $g(x)\neq 0$ et le coefficient dominant de $g(x)$ est inversible ($\in \mathbf{R}^{*}$). Alors ils existent $q,r\in\mathbf{R}[x]$ tels que 
\begin{itemize}
\item $f(x)=q(x)g(x)+r(x)$
\item $\deg r(x) < \deg g(x)$
\end{itemize}
\end{theorem}
\begin{proof}
Si $\deg f <\deg g$ alors $q(x)=0$ et $r(x)=f(x)$.

Sinon on a que $\deg f \geq \deg g$. Soient $f(x)=a_{0}+\dots+a_nx^{n}$ et $g(x)=b_{0}+\dots+b_{n}x^{n}$ avec $b_{m}\in \mathbf{R}^{*}$ et $a_{n}\neq 0$. On a que $\deg(\underbrace{ f(x)-a_{n}b_{m}^{-1}x^{n-m}g(x) }_{ h(x) })<\deg f$ donc on peut exprimer $h(x)=q'(x)g(x)+r(x)$
ou $\deg r<\deg g$. Alors
$$
f(x)-a_{n}b_{m}^{-1}x^{n-m}g(x)=q'(x)g(x)+r(x)
$$
alors $f(x)=q(x)g(x)+r(x)$ ou $q(x)=a_{n}b_{m}^{-1}x^{n-m}+q'(x)$ et $\deg r<\deg g$. 

On peut montrer que cette representation est unique. Soit $f(x)=\tilde{q}(x)g(x)+\tilde{r}(x)$ une autre representation de $f(x)$ ou $\deg \tilde{r}<\deg g$ et $q(x)\neq \tilde{q}(x)$. Alors on a 
$$
(q(x)-\tilde{q}(x))g(x)=\tilde{r}(x)-r(x)
$$
mais comme $q(x)-\tilde{q}(x)$ est un polynome non-nul et le coefficient dominant de $g(x)$ n'est pas un diviseur de zero on a que $\deg((q(x)-\tilde{q}(x))g(x))\geq \deg g$ mais $\deg(\tilde{r}(x)-r(x))<\deg g$ qui pose une contradiction. 
\end{proof}
\begin{theorem}
Soit $R$ un anneau et $f,g\in \mathbf{R}[x]$ tel que $g(x)\neq 0$ et le coefficient dominant de $g(x)$ est inversible ($\in \mathbf{R}^{*}$). Alors ils existent $q,r\in\mathbf{R}[x]$ tels que 
\begin{itemize}
\item $f(x)=g(x)q(x)+r(x)$
\item $\deg r(x) < \deg g(x)$
\end{itemize}
\end{theorem}
\begin{proof}
On reprend la preuve du theoreme precedent mais on inverse l'ordre pour obtenir $h(x)=f(x)-g(x)b_{m}^{-1}a_{n}x^{n-m}$
\end{proof}
\begin{example}
Soit $R=\mathbf{Z}^{2\times2}$ on a 
$$
f(x)=\begin{pmatrix}
1 & 0 \\
0 & 1 
\end{pmatrix}x^{4}+\begin{pmatrix}
1 & 0 \\
0 & 1
\end{pmatrix}
$$
et
$$
g(x)=\begin{pmatrix}
1 & 1 \\
0 & 1
\end{pmatrix}x^{2}+\begin{pmatrix}
1 & 0 \\
0 & 1
\end{pmatrix}
$$
Selon si on multiplie a gauche ou a droite on a deux variations:
\begin{enumerate}
\item $f(x)=q(x)g(x)+r(x)$
$$
q(x)=\begin{pmatrix}
1 & -1 \\
0 & 0 
\end{pmatrix}x^{2}+\begin{pmatrix}
-1 & 2 \\
0 & 0
\end{pmatrix},\quad r(x)=\begin{pmatrix}
2 & -2 \\
0 & 1
\end{pmatrix}
$$
\item $f(x)=g(x)q(x)+r(x)$
$$
q(x)=\begin{pmatrix}
1 & 0 \\
0 & 0
\end{pmatrix}x^{2}+\begin{pmatrix}
-1 & 0 \\
0 & 0
\end{pmatrix},\quad r(x)=\begin{pmatrix}
2 & 0 \\
0 & 1
\end{pmatrix}
$$
\end{enumerate}
\end{example}
